\stfhIndex{I039}{Aodanwenti}{奥丹问题（Audin's question）}
\section{单调拉格朗日子流形的\personAudin{}问题}
\label{sec:audin-question-monotone}

除\personGromov{}的非正合性定理外，关于 $\mathbb C^n$ 中紧拉格朗日嵌入之辛拓扑，最早的
深刻问题之一是\personAudin{} (Au88) 提出的如下问题。

\begin{question}
\label{ques:audin-minimal-maslov-two}
任意紧拉格朗日环面的最小马斯洛夫数是否都是 2？
\end{question}

\stfhIndex{I154}{Boerjieluoweiqi}{波尔捷罗维奇（Polterovich）}
\stfhIndex{I193}{Weitaibo}{维泰博（Viterbo）}
第一个肯定答案由\personViterbo{} (Vi90) 在一般情形下得到，而\personPolterovich{} (Po91a) 则在
$\mathbb R^4$ 中得到。\personViterbo{}使用作用量泛函的经典临界点理论，\personPolterovich{}使用
伪全纯曲线方法。对单调拉格朗日环面 $L\subset\mathbb C^n$，作者利用拉格朗日\personFloer{}同调的良定义性及厚--薄二分思想，在 $n\leq24$ 时肯定地回答了这个问题
(Oh96b)。同一定理也适用于一般 $(M,\omega)$ 中任意可位移单调拉格朗日环面。

\personFukaya{}、\personOh{}、\personOhta{}、\personOno{} (FOOO09) 与 L. \personBuhovsky{} (Bu10) 独立地利用谱序列与
$\mathfrak m_2$ 诱导乘积的相容性，去掉了维数限制 $n\leq24$。使用谱序列乘法性质
的想法来自 P. \personSeidel{}（私人通信，2002）。结果如下；这里的证明借自 (FOOO09)。

\begin{theorem}
\label{thm:audin-monotone-displaceable-torus}
设 $L\subset(M,\omega)$ 是可位移单调拉格朗日环面。则其最小马斯洛夫数满足
\[
 \Sigma_L=2.
\]
\end{theorem}

\begin{proof}
先指出 $T^n$ 可定向且为 spin，而单调性依定义蕴含 $\Sigma_L>0$。因为 $T^n$
可定向，非零整数 $\Sigma_L$ 必为偶数，故 $\Sigma_L\geq2$。

用反证法。反设 $\Sigma_L\geq3$。（由于 $T^n$ 可定向，这其实蕴含
$\Sigma_L\geq4$，但证明不用此事实。）单调性于是蕴含
$HF(L,L;\Lambda_{(\omega,L)})$ 良定义并在哈密顿同痕下不变。

边界映射分解为
\[
 \delta=\delta_{(0)}+\delta_{(1)}+\cdots+\delta_{(N)}
 =:\delta_{(0)}+\delta',
\]
其中 $N\leq[(n+1)/\Sigma_L]$，而 $\delta_{(k)}$ 具有形式
\begin{equation}
 \delta_{(k)}=\overline\delta_{(k)}\otimes T^{k\lambda_0}e^{k\Sigma_L}.
\label{eq:17-delta-k-monotone-torus}
\end{equation}
这里 $\overline\delta_{(k)}:E_k\to E_k$ 的次数为 $1-k\Sigma_L$。由于
$\Sigma_L\geq3$，当 $k\geq1$ 时该次数不大于 $-2$。

现在考虑 $\delta$ 在 $E_2$ 项
\[
 E_2=H^*(T^n;\mathbb Q)\otimes\Lambda_{0,\mathrm{Nov}}
\]
上的作用。回忆，上同调环 $H^*(T^n;\mathbb Q)$ 由一次上同调类
$\alpha_1,\ldots,\alpha_n\in H^1(L;\mathbb Q)$ 生成。简单的次数计算给出
\[
 \delta_2(\alpha_i)=\sum_{k=1}^{N}\overline\delta_{(k)}(\alpha_i)
 \cdot T^{k\lambda_0}e^{k\Sigma_L}=0,
 \qquad i=1,\ldots,n.
\]
\footnote{译者注：原书在本式中用同一个 $k$ 同时编号上同调生成元
$\alpha_k$ 与微分分量 $\overline\delta_{(k)}$，只写出一个匹配下标的分量；
但 $\delta_2$ 是所有正能量分量之和。译文改用 $i$ 编号生成元并补出对
$1\leq k\leq N$ 的求和；每一分量因次数至多为 $-2$ 都在 $H^1$ 上消失。}
因此
\[
 \delta_2(\mathfrak m_2(\alpha_i,\alpha_j))
 =-\mathfrak m_2(\delta_2(\alpha_i),\alpha_j)
 -\mathfrak m_2(\alpha_i,\delta_2(\alpha_j))=0.
\]
另一方面，由能量考察，
\[
 \mathfrak m_2(\alpha_i,\alpha_j)\equiv
 \alpha_i\smile\alpha_j\pmod{T^{\lambda_0}}.
\]
故对所有 $i,j=1,\ldots,n$，都有
$\delta_2(\alpha_i\smile\alpha_j)=0$。对 $\alpha_i$ 的所有可能乘积归纳应用上述
论证，得到谱序列在 $E_2$ 项退化，从而
\[
 HF^*(L;\Lambda_{0,\mathrm{Nov}})
 \cong H^*(T^n;\mathbb Q)\otimes\Lambda_{0,\mathrm{Nov}},
\]
特别地，它是 $\Lambda_{0,\mathrm{Nov}}$ 上的自由模。因此
\begin{equation}
 HF^*(L;\Lambda_{\mathrm{Nov}})
 \cong H^*(T^n;\mathbb Q)\otimes\Lambda_{\mathrm{Nov}}\neq0.
\label{eq:17-torus-floer-nonzero}
\end{equation}
（到目前为止尚未使用 $L\subset M$ 可位移的假设；它将在证明的最后一步使用。）
另一方面，由 $HF^*(L;\Lambda_{\mathrm{Nov}})$ 的不变性及存在满足
$\phi(L)\cap L=\varnothing$ 的哈密顿微分同胚 $\phi$，可知
$HF^*(L;\Lambda_{\mathrm{Nov}})=\{0\}$。这与
\eqref{eq:17-torus-floer-nonzero} 矛盾，因而 $\Sigma_L\leq2$ 得证。
\end{proof}

\begin{remark}
\label{rem:audin-modern-perspective}
从\personFloer{}理论的现代观点看，\personViterbo{}在 (Vi90) 中的证明可以重述为辛流形环路空间上
哈密顿\personFloer{}同调的一个应用 (Vi99)，而\personPolterovich{}的证明可视为拉格朗日\personFloer{}同调的应用。长期以来，人们颇感困惑：为何这两种\personFloer{}同调会给出同类结果？后来
\stfhIndex{I093}{Shenguxianzhi}{深谷贤治（Fukaya）}
\personFukaya{} (Fu06) 在拉格朗日子流形的环路空间上引入\personChas{}--\personSullivan{}括号，并把
$L\subset(M,\omega)$ 的 $A_\infty$-代数重构为德拉姆复形之 DGA 的形变，从而
在完全一般的情形下肯定地回答了\personAudin{}问题。此外还有 Y. \personEliashberg{}倡导的辛场论
方法（私人通信，2001）。据作者所知，这些证明的细节尚未发表。
\end{remark}
